\documentclass[11pt]{article}
\usepackage[a4paper,margin=1in]{geometry}
\usepackage{fontspec}
\usepackage[boldfont=false]{xeCJK}
\setCJKmainfont{FandolSong-Regular.otf}
\usepackage{amsmath}
\usepackage{amssymb}
\usepackage{array}
\usepackage{booktabs}
\usepackage{graphicx}
\usepackage{adjustbox}
\usepackage{enumitem}
\setlist[itemize]{topsep=2pt,partopsep=0pt,parsep=0pt,itemsep=2pt,leftmargin=1.4em}
\setlist[enumerate]{topsep=2pt,partopsep=0pt,parsep=0pt,itemsep=2pt,leftmargin=1.6em}
\usepackage{parskip}
\usepackage{fvextra}
\fvset{breaklines=true,breakanywhere=true,fontsize=\small}
\setlength{\parindent}{0pt}

\begin{document}

\begin{center}
  {\LARGE\bfseries Hydroxy compounds}\\[0.35em]
  {\large A Level Chemistry}
\end{center}
\vspace{0.6em}

\section*{Alcohols}

An \textbf{alcohol} 醇 has the $\text{–OH}$ (hydroxyl) functional group.

\subsection*{Making alcohols}

\begin{center}
\adjustbox{max width=\linewidth}{%
\begin{tabular}{ll}
\toprule
\textbf{Method} & \textbf{Reagents and conditions} \\
\midrule
addition of steam to an alkene & $\text{H}_2\text{O(g)}$, $\text{H}_3\text{PO}_4$ catalyst (\textbf{electrophilic addition} 亲电加成) \\
an \textbf{alkene} 烯烃 with cold dilute $\text{KMnO}_4$ & gives a \textbf{diol} 二醇 (two $\text{–OH}$ groups) \\
substitution of a \textbf{halogenoalkane} 卤代烷 & $\text{NaOH(aq)}$, heat \\
\textbf{reduction} 还原 of an \textbf{aldehyde} 醛 or \textbf{ketone} 酮 & $\text{NaBH}_4$ or $\text{LiAlH}_4$ \\
reduction of a \textbf{carboxylic acid} 羧酸 & $\text{LiAlH}_4$ \\
\textbf{hydrolysis} 水解 of an \textbf{ester} 酯 & dilute acid or alkali, heat \\
\bottomrule
\end{tabular}}
\end{center}

\subsection*{Reactions of alcohols}

\begin{itemize}
\item \textbf{combustion}: alcohols burn in oxygen to give carbon dioxide and water.

\item \textbf{substitution} to a halogenoalkane, for example with $\text{HX}$, $\text{PCl}_5$, $\text{PCl}_3$ and heat, or $\text{SOCl}_2$.

\item \textbf{with sodium}: alcohols react with sodium metal to give hydrogen and a sodium alkoxide — like water, but more slowly.

\item \textbf{oxidation} 氧化 with acidified $\text{K}_2\text{Cr}_2\text{O}_7$ (or $\text{KMnO}_4$). The product depends on the class of alcohol (see below).

\item \textbf{dehydration} 脱水 to an alkene, using a hot $\text{Al}_2\text{O}_3$ catalyst or concentrated acid.

\item \textbf{ester formation}: an alcohol reacts with a carboxylic acid (with concentrated $\text{H}_2\text{SO}_4$ catalyst) to make an ester.

\end{itemize}

\subsection*{Three classes and how oxidation tells them apart}

An alcohol is \textbf{primary} 伯, \textbf{secondary} 仲 or \textbf{tertiary} 叔, depending on how many carbons are joined to the carbon holding the $\text{–OH}$. Some molecules have more than one $\text{–OH}$ group.

\begin{center}
\adjustbox{max width=\linewidth}{%
\begin{tabular}{ll}
\toprule
\textbf{Class} & \textbf{Oxidation product} \\
\midrule
primary & \textbf{aldehyde} (by \textbf{distillation} 蒸馏), then \textbf{carboxylic acid} (by \textbf{reflux} 回流) \\
secondary & a \textbf{ketone} \\
tertiary & not oxidised \\
\bottomrule
\end{tabular}}
\end{center}

In a quick test, acidified $\text{K}_2\text{Cr}_2\text{O}_7$ turns from \textbf{orange to green} with a primary or secondary alcohol, but stays orange with a tertiary alcohol.

\begin{itemize}
\item \textbf{distillation} removes the aldehyde as it forms, before it can be oxidised further.

\item \textbf{reflux} keeps boiling the mixture and returning the vapour, so the alcohol is fully oxidised to the carboxylic acid.

\end{itemize}

\subsection*{The iodoform test}

If you warm an alcohol that contains the $\text{CH}_3\text{CH(OH)}-$ group with alkaline aqueous iodine, you get a pale yellow precipitate of \textbf{tri-iodomethane} 三碘甲烷 ($\text{CHI}_3$) and the ion $\text{RCO}_2^-$. This is a useful test for that group.

\subsection*{Acidity of alcohols}

The $\text{–OH}$ group makes alcohols very weakly acidic: they can lose the $\text{H}^+$ to form an $\text{RO}^-$ ion. But their \textbf{acidity} 酸性 is \textbf{lower} than that of water. This is because the alkyl group pushes electron density onto the oxygen, which makes the $\text{RO}^-$ ion less stable, so the alcohol holds onto its $\text{H}^+$ more tightly.


\end{document}
